% Chapter 1 — includable by main.tex
% Title converted to a chapter heading
\chapter{Jezici i osnovne operacije}
\label{chap:1}

\begin{quote}
\itshape
Sve što ćemo raditi u ovom kolegiju vrti se oko jednog jednostavnog pitanja:
pripada li neka riječ nekom jeziku? Ovo poglavlje postavlja temelje ---
definira što su uopće ,,riječi'' i ,,jezici'', i uvodi osnovne operacije s
kojima ćemo ih graditi i kombinirati.
\end{quote}


% ================================================================
\section{Tri pogleda na isti objekt}
% ================================================================

Zamislite funkciju $f : \mathbb{R} \to \mathbb{R}$. Možemo je opisati na
barem tri načina: kao skup uređenih parova $\{(x, f(x)) \mid x \in \mathbb{R}\}$,
kao formulu poput $f(x) = x^2 + 1$, ili kao program koji za svaki ulaz $x$
vraća odgovarajući izlaz. Sva tri opisa govore o istoj stvari, ali svaki naglašava
nešto drugačije i koristan je u različitim situacijama.

Formalni jezici bave se istim trikom, ali za jezike umjesto funkcija.
\textbf{Jezik} je skup --- statični popis riječi koje mu pripadaju.
\textbf{Gramatika} je dinamično pravilo koje taj skup \emph{generira},
kreće od početnog simbola i primjenjuje pravila prepisivanja.
\textbf{Automat} je idealizirani stroj koji \emph{prepoznaje} jezik čitajući
ulaznu riječ znak po znak i na kraju odlučuje: ,,da'' ili ,,ne''.

Ono što čini ovaj predmet zanimljivim jest to da između ova tri pogleda
postoje duboke i neočite veze. Na primjer, nije odmah jasno zašto bi za svaku
gramatiku određenog tipa \emph{nužno} postojao automat koji prepoznaje isti
jezik --- ali to je upravo ono što ćemo dokazivati.

Sve klase koje ćemo proučavati organizirane su u \textbf{Chomskyjevu hijerarhiju},
kojoj se vraćamo na kraju poglavlja.

% ================================================================
\section{Abecede i riječi}
% ================================================================

Prije nego što možemo govoriti o jezicima, moramo dogovoriti s kojim simbolima
radimo.

\begin{definicija}[Abeceda]
\textbf{Abeceda} je konačni, neprazni skup $\Sig$. Elemente abecede zovemo
\textbf{znakovima} i često označavamo malim grčkim slovima: $\alpha, \beta,
\gamma, \ldots$
\end{definicija}

Ništa posebno --- abeceda je samo dogovor o tome koji su nam simboli na
raspolaganju. Ograničenje da mora biti \emph{konačna} je ključno: svi naši
strojevi i gramatike imat će konačni opis, pa moramo krenuti od konačne
abecede.

\begin{primjer}[Nekoliko uobičajenih abeceda]
U praksi ćemo najčešće raditi s:
\begin{itemize}
  \item $\Sig = \{0, 1\}$ --- binarna abeceda, dovoljna za kodiranje čega god želimo
  \item $\Sig = \{a, b\}$ ili $\Sig = \{a, b, c\}$ --- mali primjeri za ilustracije
  \item $\Sig = \{0, 1, \ldots, 9\}$ --- decimalne znamenke
\end{itemize}
Abeceda ne mora biti skup slova ili znamenki --- može biti skup tokena, boja,
instrukcija ili bilo čega drugog. Bitno je samo da je konačna.
\end{primjer}

\begin{definicija}[Riječ i njezina duljina]
\textbf{Riječ} nad abecedom $\Sig$ je konačni niz znakova iz $\Sig$:
\[
  w = \alpha_1 \alpha_2 \cdots \alpha_n, \qquad \alpha_i \in \Sig.
\]
\textbf{Duljina} riječi $w$, pisana $|w|$, jest broj znakova u njoj.
Pišemo također $|w|_\alpha$ za broj pojavljivanja znaka $\alpha$ u $w$.
\end{definicija}

Posebno mjesto zauzima \textbf{prazna riječ} $\eps$ --- jedina riječ duljine 0.
To nije znak, nego prazan niz znakova, i vrijedi $|\eps| = 0$.

\begin{napomena}
Prazna riječ $\eps$ i prazan \emph{jezik} $\emptyset$ nisu isto.
Prazna riječ \emph{postoji} kao objekt, samo ne sadrži nijedan znak.
Prazan jezik je skup koji ne sadrži nijednu riječ --- pa onda ni $\eps$.
Ovu razliku je lako zamijeniti na ispitu, pazite.
\end{napomena}

Uobičajeno je poistovjećivati znak $\alpha \in \Sig$ s jednakovrijednom
jednoznakovnom riječi. Dakle, $a$ kao znak i $a$ kao jednoelementna
,,rečenica'' su isti objekt --- samo gledamo iz drugog kuta.

% ================================================================
\section{Konkatenacija}
% ================================================================

Osnovna operacija kojom iz kratkih riječi gradimo dulje jest
\textbf{konkatenacija} --- spajanje dvije riječi jedne za drugom.

\begin{definicija}[Konkatenacija]
Za riječi $u = \alpha_1 \cdots \alpha_n$ i $v = \beta_1 \cdots \beta_m$
definiramo
\[
  uv \;:=\; \alpha_1 \cdots \alpha_n \, \beta_1 \cdots \beta_m.
\]
Vrijedi $|uv| = |u| + |v|$ --- duljine se zbrajaju.
\end{definicija}

Konkatenacija je asocijativna: $(uv)w = u(vw)$, pa zagrade su suvišne
i pišemo jednostavno $uvw$. Prazna riječ $\eps$ je neutralni element:
$\eps w = w \eps = w$ --- dodavanje ničega ništa ne mijenja.

Budući da konkatenacija ,,izgleda kao množenje'', koristimo uobičajenu
multiplikativnu notaciju za potencije:
\[
  u^0 := \eps, \qquad u^1 := u, \qquad u^2 := uu, \qquad u^3 := uuu, \quad\ldots
\]

\begin{primjer}
Neka je $u = ab$ i $v = ba$. Tada $uv = abba$, $vu = baab$, $u^3 = ababab$.
Konkatenacija nije komutativna --- $uv \neq vu$ u općem slučaju.
\end{primjer}

% ================================================================
\section{Jezici}
% ================================================================

Sad kad znamo što su riječi, možemo definirati jezik.

\begin{definicija}[$\Sig^*$ i jezik]
Skup \emph{svih} mogućih riječi nad abecedom $\Sig$ (uključujući $\eps$)
označavamo s $\Sig^*$:
\[
  \Sig^* \;:=\; \{\eps\} \cup \Sig \cup \Sig^2 \cup \Sig^3 \cup \cdots
\]
\textbf{Jezik} nad $\Sig$ je bilo koji podskup $L \subseteq \Sig^*$.
\end{definicija}

Definicija je namjerno vrlo široka --- jezik je doslovno \emph{bilo koji}
skup riječi, konačni ili beskonačni. Jedino ograničenje je da su sve riječi
nad istom abecedom. Nema uvjeta da jezik mora imati neku posebnu strukturu;
to tek postaje zanimljivo kad počnemo klasificirati jezike po tome
\emph{kakvu} strukturu imaju.

\begin{primjer}
Nad $\Sig = \{a, b\}$, evo nekoliko jezika raznih vrsta:
\begin{itemize}
  \item $L_1 = \{a, ab, aba\}$ --- konačni jezik s tri elementa
  \item $L_2 = \{w \in \Sig^* \mid |w|_a = |w|_b\}$ --- sve riječi s jednakim
        brojem $a$-ova i $b$-ova; beskonačni
  \item $L_3 = \{a^n b^n \mid n \geq 1\} = \{ab, aabb, aaabbb, \ldots\}$ ---
        klasični primjer koji ćemo često susretati
  \item $\emptyset$ --- prazan jezik; nema nijedne riječi
  \item $\{\eps\}$ --- jezik s jednom riječi, praznom
  \item $\Sig^*$ --- svi mogući nizovi znakova, bez ikakva ograničenja
\end{itemize}
\end{primjer}

Jezike možemo kombinirati skupovnim operacijama ($\cup$, $\cap$,
$\setminus$, komplement spram $\Sig^*$), ali i operacijom konkatenacije,
proširenom s riječi na skupove:

\begin{definicija}[Konkatenacija jezika]
Za jezike $L, M \subseteq \Sig^*$ definiramo
\[
  LM \;:=\; \{uv \mid u \in L,\; v \in M\}.
\]
Dakle, $LM$ se dobiva lijepljenjem svake riječi iz $L$ sa svakom riječi iz $M$.
\end{definicija}

\begin{primjer}
Neka je $L = \{a, ab\}$ i $M = \{b, c\}$. Sve kombinacije su:
\[
  LM = \{a \cdot b,\; a \cdot c,\; ab \cdot b,\; ab \cdot c\}
     = \{ab,\; ac,\; abb,\; abc\}.
\]
A što je $L^2 = LL$? Svaka riječ iz $L$ zalijepljena sa svakom:
\[
  L^2 = \{aa,\; aab,\; aba,\; abab\}.
\]
Primijetite da $L^2$ ima točno $|L|^2 = 4$ elementa --- u ovom slučaju
nema ,,sudara'' (različitih kombinacija koje daju istu riječ). Nije
uvijek tako.
\end{primjer}

% ================================================================
\section{Kleeneve operacije}
% ================================================================

Najvažnija operacija u cijelom kolegiju je Kleeneva zvijezda.
Susresti ćete je svugdje --- u regularnim izrazima, gramatikama, posvuda.

\begin{definicija}[Kleenejev plus i zvijezda]
Za jezik $L \subseteq \Sig^*$ definiramo:
\begin{align*}
  L^+ &\;:=\; \bigcup_{k \geq 1} L^k
        \;=\; L \cup L^2 \cup L^3 \cup \cdots \\[6pt]
  L^* &\;:=\; \bigcup_{k \geq 0} L^k
        \;=\; \{\eps\} \cup L \cup L^2 \cup L^3 \cup \cdots
\end{align*}
Razlika je samo u tome što $L^*$ uključuje $L^0 = \{\eps\}$, a $L^+$ ne.
Vrijedi $L^* = \{\eps\} \cup L^+$, tj.\ zvijezda je plus uz dodatak
prazne riječi.
\end{definicija}

Kako čitati ove definicije? $L^*$ je skup svih \emph{nizanja} nula ili više
riječi iz $L$ jedne za drugom. Kad kažemo da neka riječ $w$ pripada $L^*$,
mislimo da se $w$ može ,,izrezati'' na komadiće pri čemu svaki komadić
pripada $L$.

\begin{kljucno}
$w \in L^*$ \textbf{ako i samo ako} postoji rastav
$w = w_1 w_2 \cdots w_k$ (za neko $k \geq 0$) takav da $w_i \in L$
za svaki $i$. Slučaj $k = 0$ daje $w = \eps$ --- uvijek prihvatljivo.
\end{kljucno}

\begin{primjer}
Uzmimo $L = \{0, 11\}$ nad $\Sig = \{0, 1\}$. Tada:
\[
  L^* = \{\eps,\; 0,\; 11,\; 00,\; 011,\; 110,\; 1111,\; 000,\; \ldots\}
\]
Je li $110011 \in L^*$? Da --- rastav je
$\underbrace{11}_{}\,\underbrace{0}_{}\,\underbrace{0}_{}\,\underbrace{11}_{}$,
svaki komadić je u $L$.

Je li $1 \in L^*$? Ne. Jedina jedinka u $L$ je $0$; jedini dvoznakac je $11$;
nema načina sastaviti samu $1$ od komadića iz $\{0, 11\}$.
\end{primjer}

\begin{primjer}
Posebno važan slučaj: ako je $\Sig$ abeceda (shvaćena kao jezik
jednoznakovnih riječi), tada je $\Sig^*$ upravo \emph{skup svih mogućih
riječi} nad $\Sig$ --- što je i definicija iz prethodnog odjeljka.
Zvijezda u oznaci $\Sig^*$ nije slučajnost.
\end{primjer}

Postoji korisna analogija s aritmetikom. Konkatenacija je množenje, prazna
riječ $\eps$ je jedinica, a Kleeneva zvijezda nalikuje na ,,beskonačnu
geometrijsku sumu'' $1 + L + L^2 + \cdots$, samo u algebri jezika.
Ta analogija nije samo poetska --- na njoj se temelji teorija regularnih
izraza koju ćemo graditi u sljedećem poglavlju.

\begin{zadatak}
\begin{enumerate}[label=(\alph*), itemsep=4pt]
  \item Izrazite $L^*$ pomoću $L^+$ i obrnuto. Koji uvjet mora biti ispunjen
        da vrijedi $L^+ = L^* \setminus \{\eps\}$?
  \item Dokažite ili opovrgnite: $(L \cup M)^* = (L^* M^*)^*$.
  \item Pokažite da je $\emptyset^* = \{\eps\}$.
        Intuitivno, zašto to ima smisla?
  \item Za $L = \{a\}$, opišite $L^*$ riječima. Koji je to dobro poznati skup?
\end{enumerate}
\end{zadatak}

% ================================================================
\section{Pozitivni jezici}
% ================================================================

Jedno naizgled sitno, ali tehnički korisno svojstvo jest pozitivnost.

\begin{definicija}[Pozitivan jezik]
Kažemo da je jezik $L$ \textbf{pozitivan} ako $\eps \notin L$.
\end{definicija}

Zašto nam je to uopće važno? U dokazima koji idu indukcijom po duljini
riječi, pozitivnost osigurava da konkatenacijom ne možemo ostati na istoj
duljini. Precizno: ako je $L$ pozitivan i $w \in LM$, dio koji pripada $M$
strogo je kraći od $w$ --- pa indukcija ,,napreduje''.

Vrijedi nekoliko jednostavnih pravila:
\begin{itemize}[itemsep=3pt]
  \item $\emptyset$ je pozitivan (trivijalno --- nema elemenata, pa ni $\eps$).
  \item Svaki jednoznakovni jezik $\{\alpha\}$ je pozitivan.
  \item $L^*$ \emph{nikad nije} pozitivan --- uvijek vrijedi $\eps \in L^*$.
  \item $LM$ je pozitivan ako i samo ako je barem jedan od $L$, $M$ pozitivan.
  \item $L \cup M$ je pozitivan ako i samo ako su \emph{oba} $L$ i $M$ pozitivna.
\end{itemize}

Zadnja dva pravila su dobar test razumijevanja --- probajte ih dokazati sami
iz definicija.

% ================================================================
\section{Chomskyjeva hijerarhija}
% ================================================================

Na kraju poglavlja smjestimo sve u širi okvir. Chomskyjeva hijerarhija
klasifikacija je jezika u klase prema tome koliko su ,,teški za prepoznati''.

Neformalna ideja: što je jezik strukturiraniji i zahtjevniji, to je moćniji
stroj (ili gramatika) potreban za njegovo opisivanje. Najjednostavniji jezici
--- \textbf{regularni} --- mogu se prepoznati \emph{bez ikakva pamćenja} osim
trenutnog stanja. Za malo zahtjevnije --- \textbf{bezkontekstne} --- treba nam
stog (neograničena memorija, ali pristup samo s vrha). Zatim dolaze
\textbf{kontekstni} jezici koji trebaju traku jednake duljine kao ulaz.
Na samom vrhu su \textbf{rekurzivno prebrojivi} jezici, za koje treba
puni Turingov stroj.

Svakoj klasi odgovara vrsta gramatike i vrsta automata:

\begin{center}
\renewcommand{\arraystretch}{1.5}
\begin{tabular}{llll}
\hline
\textbf{Klasa} & \textbf{Oznaka} & \textbf{Gramatika} & \textbf{Automat} \\
\hline
Regularni        & Reg & desno/lijevolinearna & konačni (KA, NKA) \\
Bezkontekstni    & BK  & bezkontekstna (BKG)  & potisni (PA) \\
Kontekstni       & Kon & kontekstna (KG)      & ograničeni (OA) \\
Rekurzivni       & Rek & ---                  & Turingov odlučitelj \\
Rek.\ prebrojivi & RE  & opća (OG)            & Turingov prepoznavač \\
\hline
\end{tabular}
\end{center}

Među klasama vrijede \emph{prave inkluzije} --- svaka klasa je pravi
podskup sljedeće:
\[
  \mathrm{Reg} \;\subsetneq\; \mathrm{BK} \;\subsetneq\; \mathrm{Kon}
  \;\subsetneq\; \mathrm{Rek} \;\subsetneq\; \mathrm{RE}.
\]
,,Prava'' inkluzija znači da za svaki par susjednih klasa postoji jezik koji
pripada gornjoj, ali ne i donjoj. Klasični primjeri koji ,,razdvajaju''
klase su:
\begin{align*}
  L_{=}      &= \{a^n b^n \mid n \geq 1\} \;\in\; \mathrm{BK} \setminus \mathrm{Reg}, \\
  L_{\equiv} &= \{a^n b^n c^n \mid n \geq 1\} \;\in\; \mathrm{Kon} \setminus \mathrm{BK}.
\end{align*}

Zašto, primjerice, $L_=$ nije regularan? Intuitivno: kad konačni automat
čita $a^n b^n$, u trenutku kad dođe do prvih $b$-ova mora nekako ,,znati''
koliko je $a$-ova pročitao --- a ima samo konačno mnogo stanja, pa ne može
pamtiti proizvoljan broj. Formalni dokaz dat ćemo u poglavlju o lemi o
pumpanju.

\begin{kljucno}[Što nas čeka]
Kolegij prolazi kroz hijerarhiju odozdo prema gore: regularni jezici
(poglavlja 2--4), bezkontekstni (poglavlje 5), kontekstni (poglavlje 7),
i na kraju Turingovi strojevi i odlučivost (poglavlje 8).
Svaki put kad pređemo u višu klasu, pronaći ćemo jezik koji ne stane
u onu ispod --- i naučiti \emph{zašto}.
\end{kljucno}

% ================================================================
\section*{Sažetak}
% ================================================================

Ovim smo postavili temelje. U jednoj rečenici: jezik je skup riječi nad
nekom abecedom, a ovo poglavlje definira kako se s takvim skupovima radi.

\medskip
\noindent
\begin{tabular}{ll}
  $\Sig$ & abeceda --- konačni, neprazni skup znakova \\
  $w, \eps$ & riječ, prazna riječ \\
  $|w|$ & duljina riječi \\
  $\Sig^*$ & skup svih riječi nad $\Sig$ \\
  $L \subseteq \Sig^*$ & jezik \\
  $LM$ & konkatenacija jezika \\
  $L^*$, $L^+$ & Kleeneva zvijezda i plus \\
\end{tabular}

\medskip
\noindent U sljedećem poglavlju uvodimo \textbf{regularne izraze} ---
kompaktan formalizam za opisivanje regularnih jezika koji koristi upravo
operacije $\cup$, konkatenaciju i ${}^*$ koje smo ovdje upoznali.
